В ходе работы над проектом под названием <<Разработка видеоигры"=головоломки “Одноплатник”>> была успешно реализована обучающая 2D"=игра, направленная на знакомство студентов-первокурсников с базовыми принципами подключения периферийных устройств к одноплатному компьютеру Raspberry Pi. 

Игра наглядно демонстрирует, как компьютерные игры могут использоваться в обучающих целях. Вместо заучивания схем или инструкций, пользователь получает возможность экспериментов: он может перетаскивать устройства, вращать провода, получать немедленную визуальную и аудиообратную связь. Всё это повышает мотивацию, интерес и улучшает усвоение материала.

Таким образом, проект <<Одноплатник>> стал не только рабочим прототипом образовательной игры, но и практическим подтверждением того, что игровые технологии могут эффективно использоваться для введения в технические дисциплины. Благодаря сочетанию учебной цели и технической реализации, проект представляет ценность для образовательного процесса на факультете КНиИТ.